% ============================================================
% Front matter: cover, title page, declaration, approval,
% acknowledgements, abstracts, lists
% ============================================================

% Team, used on the cover, the title page and the declaration
\newcommand{\teamlist}{%
  Khaled Ibrahim Alghmlas -- 445050235\\
  Faleh Hamoud Alsubaie -- 445050053\\
  Ibrahim Abdulrahman Alharthi -- 445050024}

\newcommand{\titleblock}[1]{% #1 = vertical gap between blocks
  {\LARGE\bfseries\setstretch{1.15}\projecttitle\par}
  \vspace{#1}
  {\large\bfseries\setstretch{1.2}\teamlist\par}
  \vspace{#1}
  To obtain\par
  {\large BSc in Computer Science\par}
  \vspace{0.7\dimexpr#1\relax}
  Department of Computer Sciences\par
  College of Computer Engineering and Sciences\par
  \vspace{0.7\dimexpr#1\relax}
  Supervisor: Dr.\ Khalil Chebil\par
  \vspace{0.7\dimexpr#1\relax}
  December 2026\par}

% ------------------------------------------------------------
% Cover page
% ------------------------------------------------------------
\begin{titlepage}
\setstretch{1.1}
\noindent
\begin{minipage}[b]{0.4\textwidth}
\footnotesize
Kingdom of Saudi Arabia\\
Ministry of Education\\
Prince Sattam Bin Abdulaziz University\\
College of Computer Engineering and Sciences\\
Department of Computer Sciences
\end{minipage}\hfill
\begin{minipage}[b]{0.2\textwidth}
\centering\includegraphics[width=0.75\linewidth]{psau-logo.jpeg}
\end{minipage}\hfill
\begin{minipage}[b]{0.36\textwidth}
\raggedleft\footnotesize
\ar{المملكة العربية السعودية}\\
\ar{وزارة التعليم}\\
\ar{جامعة الأمير سطام بن عبد العزيز}\\
\ar{كلية هندسة وعلوم الحاسب}\\
\ar{قسم علوم الحاسب}
\end{minipage}

\vspace{1.6cm}
\begin{center}
\titleblock{1.2cm}
\vspace{0.6cm}
Number: *****
\end{center}
\end{titlepage}

% ------------------------------------------------------------
% Title page
% ------------------------------------------------------------
\begin{titlepage}
\vspace*{2.5cm}
\begin{center}
\titleblock{1.8cm}
\end{center}
\end{titlepage}

% ------------------------------------------------------------
% Declaration
% ------------------------------------------------------------
\setcounter{page}{2}
\frontchapter{Declaration}

I/we hereby declare that this project report is based on my/our original
work except for citations and quotations which have been duly acknowledged.
I/We also declare that it has not been previously or concurrently submitted
for any other degree at PSAU or any other institution.

\vspace{1.5cm}
\noindent Signature:\ \dotfill\ \ \ Date:\ \makebox[5cm]{\dotfill}

\vspace{1cm}
\noindent
\begin{tabular}{@{}p{0.5\textwidth}l@{}}
Name: Khaled Ibrahim Alghmlas      & ID: 445050235 \\[0.4em]
Name: Faleh Hamoud Alsubaie        & ID: 445050053 \\[0.4em]
Name: Ibrahim Abdulrahman Alharthi & ID: 445050024 \\
\end{tabular}

% ------------------------------------------------------------
% Approval
% ------------------------------------------------------------
\frontchapter{Approval for Submission}

I certify that this project report entitled \textbf{``\projecttitle''} was
prepared by Khaled Ibrahim Alghmlas, Faleh Hamoud Alsubaie and Ibrahim
Abdulrahman Alharthi and has met the required standard for submission in
partial fulfillment of the requirements for the degree of Bachelor of Science
in Computer Science at PSAU.

\vspace{1.2cm}
Approved by

\vspace{2cm}
\noindent Signature:\ \makebox[7cm]{\dotfill}

\noindent Supervisor: Dr.\ Khalil Chebil

\noindent Date:\ \makebox[7cm]{\dotfill}

% ------------------------------------------------------------
% Acknowledgements
% ------------------------------------------------------------
\frontchapter{Acknowledgements}

We would like to express our sincere gratitude to our supervisor,
Dr.\ Khalil Chebil, for his continuous guidance, encouragement, and valuable
feedback throughout the course of this project, including his encouragement
to begin early development work ahead of schedule and the detailed
improvement plan he prepared for the system. We are also grateful to the
faculty members of the Department of Computer Sciences at Prince Sattam Bin
Abdulaziz University for the knowledge and support that made this project
possible.

We would like to thank each other as a team: Khaled Ibrahim Alghmlas, Faleh
Hamoud Alsubaie and Ibrahim Abdulrahman Alharthi, for the collaboration and
shared effort that went into every part of NetMind. We would also like to
thank our fellow students working on other graduation projects who generously
shared their insights and discussed ideas with us on topics related to this
work.

% ------------------------------------------------------------
% Abstract (English)
% ------------------------------------------------------------
\frontchapter{Abstract}

\begin{center}
\textbf{\projecttitle}
\end{center}

Modern networks combine many types of devices --- routers, switches,
firewalls, and IoT equipment such as IP cameras --- and administering them
typically requires familiarity with several vendor-specific command-line
interfaces, which slows operations and increases the risk of operator error.
This project presents NetMind, a network operations assistant that lets a
user manage and monitor a network through natural-language conversation in
Arabic or English rather than memorised commands. NetMind is built around a
Large Language Model, accessed through the Groq API, given callable tools to
check device status and IP addresses, power devices on or off, list devices
on a network segment, test connectivity between devices, audit a firewall's
configuration for common weaknesses, and verify whether a security camera
still uses its default password. Destructive actions are never executed on
the model's word alone: they require explicit confirmation by the user in the
interface and are recorded in an audit log.

To allow safe experimentation, NetMind operates against a simulated network
built with Containerlab, with real IP addressing and routing between
branches, so every action is executed against a genuine, contained network
rather than a mock. A second, read-only mode extends the same interface to a
real physical network, discovering devices by combining ARP scanning, a ping
sweep, IPv6 neighbour discovery and mDNS without altering their state. A
working prototype of both modes, together with a live monitoring dashboard,
alert notifications, persistent storage and an automated test suite, was
built and validated ahead of schedule. This document covers the motivation,
requirements, design, and architecture of NetMind, the reasoning behind the
tools and techniques chosen to build it, and the plan for the second term.

% ------------------------------------------------------------
% Abstract (Arabic)
% ------------------------------------------------------------
\chapter*{\hfill\ar{المستخلص}}
\addcontentsline{toc}{chapter}{ABSTRACT (ARABIC)}

\begin{center}
\ar{\textbf{نتمايند: مساعد قائم على نماذج اللغة الكبيرة لإدارة وتأمين الشبكات}}
\end{center}

\begin{arabicpar}
تجمع الشبكات الحديثة أنواعًا عديدة من الأجهزة، مثل الموجهات والمحولات
(السويتشات) وجدران الحماية ومعدات إنترنت الأشياء كالكاميرات الأمنية،
وعادةً ما تتطلب إدارتها الإلمام بعدة واجهات أوامر نصية خاصة بكل مورّد، مما
يُبطئ عمليات التشغيل ويزيد من احتمال وقوع أخطاء بشرية. يقدّم هذا المشروع
نظام نتمايند (NetMind)، وهو مساعد لعمليات الشبكات يتيح للمستخدم إدارة
الشبكة ومراقبتها من خلال محادثة بلغة طبيعية باللغتين العربية والإنجليزية،
بدلاً من حفظ الأوامر. يعتمد نتمايند على نموذج لغوي كبير، يُستدعى عبر واجهة
Groq البرمجية، زُوّد بمجموعة من الأدوات القابلة للاستدعاء للتحقق من حالة
الأجهزة وعناوين IP الخاصة بها، وتشغيلها أو إيقافها، وسرد الأجهزة على جزء
معيّن من الشبكة، واختبار الاتصال بين الأجهزة، وتدقيق إعدادات جدار الحماية
للكشف عن الثغرات الشائعة، والتحقق مما إذا كانت كاميرا المراقبة ما زالت
تستخدم كلمة المرور الافتراضية. ولا تُنفَّذ الإجراءات التدميرية بناءً على
كلام النموذج وحده، بل تتطلب تأكيدًا صريحًا من المستخدم في الواجهة وتُسجَّل
في سجل تدقيق.
\end{arabicpar}

\begin{arabicpar}
ولإتاحة التجربة الآمنة، يعمل نتمايند على شبكة محاكاة مبنية باستخدام
Containerlab، بعنونة IP وتوجيه حقيقيين بين الفروع، بحيث يُنفَّذ كل إجراء
على شبكة حقيقية ومحتواة فعلاً وليس على نموذج وهمي. كما يوفر النظام وضعًا
ثانيًا للقراءة فقط يمتد بالواجهة نفسها إلى شبكة فعلية حقيقية، حيث يكتشف
الأجهزة المتصلة بالجمع بين فحص ARP وفحص ping واكتشاف جيران IPv6 وبروتوكول
mDNS دون التأثير في حالتها. وقد تم بناء نموذج أولي يُطبّق كلا الوضعين، إلى
جانب لوحة مراقبة حيّة وتنبيهات وتخزين دائم ومجموعة اختبارات آلية، والتحقق من
عمله قبل الموعد المقرر. ويغطي هذا التقرير دوافع المشروع ومتطلباته وتصميمه
ومعماريته، والأسباب وراء اختيار الأدوات والتقنيات المستخدمة في بنائه، وخطة
العمل للفصل الدراسي الثاني.
\end{arabicpar}

% ------------------------------------------------------------
% Lists
% ------------------------------------------------------------
\cleardoublepage
\tableofcontents
\cleardoublepage
\listoftables
\cleardoublepage
\listoffigures

% ------------------------------------------------------------
% Abbreviations
% ------------------------------------------------------------
\frontchapter{List of Symbols/Abbreviations}

\setlength{\LTleft}{0pt}
\begin{longtable}{@{}p{3.2cm}l@{}}
AI      & Artificial Intelligence \\
API     & Application Programming Interface \\
ARP     & Address Resolution Protocol \\
CI      & Continuous Integration \\
CPU     & Central Processing Unit \\
CVE     & Common Vulnerabilities and Exposures \\
DHCP    & Dynamic Host Configuration Protocol \\
DNS     & Domain Name System \\
FR      & Functional Requirement \\
HTTP    & Hypertext Transfer Protocol \\
ICMP    & Internet Control Message Protocol \\
IoT     & Internet of Things \\
IP      & Internet Protocol \\
JSON    & JavaScript Object Notation \\
LLM     & Large Language Model \\
MAC     & Media Access Control \\
MCP     & Model Context Protocol \\
mDNS    & Multicast Domain Name System \\
NDP     & Neighbor Discovery Protocol (IPv6) \\
NETCONF & Network Configuration Protocol \\
NFR     & Non-Functional Requirement \\
OUI     & Organizationally Unique Identifier \\
OWASP   & Open Worldwide Application Security Project \\
PSAU    & Prince Sattam Bin Abdulaziz University \\
RBAC    & Role-Based Access Control \\
REST    & Representational State Transfer \\
SDK     & Software Development Kit \\
SNMP    & Simple Network Management Protocol \\
SQL     & Structured Query Language \\
SSDP    & Simple Service Discovery Protocol \\
SSH     & Secure Shell \\
TCP     & Transmission Control Protocol \\
UDP     & User Datagram Protocol \\
UPnP    & Universal Plug and Play \\
VM      & Virtual Machine \\
WSL     & Windows Subsystem for Linux \\
YAML    & YAML Ain't Markup Language \\
\end{longtable}

% ------------------------------------------------------------
% Keywords
% ------------------------------------------------------------
\frontchapter{List of Keywords}

NetMind, Large Language Models, Network Automation, Network Security,
Containerlab, Firewall Auditing, Network Simulation, Natural Language
Interface, Docker, Network Operations, Tool Calling, AI Agent, Network
Discovery, Network Topology, Python, Scapy, SQLite.
